\documentclass[11pt,a4paper]{article}

%% PACHETE MATEMATICE
\usepackage{amsmath,amssymb,amsfonts}
\usepackage{amsthm}
\usepackage{mathpazo}
\usepackage{eulervm}
\usepackage{enumerate}
\usepackage{color}
\usepackage{graphicx}
\usepackage[all]{xy}
\usepackage{xcolor}
\usepackage{framed}
\usepackage[unicode]{hyperref}
\setcounter{MaxMatrixCols}{20}
\usepackage{multicol}
% \usepackage[romanian]{babel}


%% ASPECT
\linespread{1.05}
\usepackage[T1]{fontenc}
\usepackage[utf8x]{inputenc}
\usepackage[margin=2cm]{geometry}
% \usepackage[color]{showkeys}
\usepackage{anyfontsize}
\usepackage{titlesec}
\titleformat*{\section}{\large\bfseries}

\usepackage{fancyhdr}
\pagestyle{fancy}
\rhead{\small \color{gray} Notițe de Adrian Manea}
\lhead{\small \color{gray} ETTI-AM2, 2017---2018, M. Joița \& A. Niță}


\usepackage[autostyle = false, style = german]{csquotes}
\usepackage[romanian]{babel}
\newcommand\qq{\enquote}

%% COMENZILE MELE
\renewcommand*{\proofname}{Dem.:}
\newcommand\bb{\mathbb}
\newcommand\dr{\mathrm}
\newcommand\kal{\mathcal}
\newcommand\fk{\mathfrak}
\newcommand\op{\oplus}
\newcommand\ot{\otimes}
\newcommand\ds{\displaystyle}
\newcommand\id{\indent}
\newcommand\nid{\noindent}
\newcommand\seq{\subseteq}
\newcommand\sneq{\subsetneq}
\newcommand\speq{\supseteq}
\newcommand\spneq{\supsetneq}
\newcommand\rar{\rightarrow}
\newcommand\Rar{\Rightarrow}
\newcommand\xrar{\xrightarrow}
\newcommand\lar{\leftarrow}
\newcommand\Lar{\Leftarrow}
\newcommand\xlar{\xleftarrow}
\newcommand\lrar{\leftrightarrow}
\newcommand\Lrar{\Leftrightarrow}
\newcommand\ideal{\trianglelefteq}
\newcommand\ai{\text{ a.\^{i}. }}
\newcommand\sk{(\textit{Schi\c{t}\u{a}}) }
\newcommand\rhup{\rightharpoonup}
\newcommand\rhdn{\rightharpoondown}
\newcommand\lhup{\leftharpoonup}
\newcommand\lhdn{\leftharpoondown}
\newcommand\vid{\emptyset}
\newcommand\wh{\widehat}
\newcommand\ol{\overline}
\newcommand\NN{\mathbb{N}}
\newcommand\RR{\mathbb{R}}
\newcommand\ZZ{\mathbb{Z}}
\newcommand\QQ{\mathbb{Q}}
\newcommand\CC{\mathbb{C}}

%% STILURILE MELE
\newtheoremstyle{dotless}{}{}{\itshape}{}{\bfseries}{:}{ }{}

\theoremstyle{dotless}
\newtheorem{theorem}{Teorem\u{a}}[section]
\newtheorem{proposition}{Propozi\c{t}ie}[section]
\newtheorem{lemma}{Lem\u{a}}[section]
\newtheorem{corollary}{Corolar}[section]
\newtheorem{exercise}{Exerci\c{t}iu}

\newtheoremstyle{dotlessdr}{}{}{}{}{\bfseries}{:}{ }{}
\theoremstyle{dotlessdr}
\newtheorem{example}{Exemplu}[section]
\newtheorem{definition}{Defini\c{t}ie}[section]
\newtheorem{remark}{Observa\c{t}ie}[section]




\begin{document}
% \definecolor{labelkey}{RGB}{222,0,0}

\begin{center}
  {\large\textbf{Seminar 5 \\
      Ecuații cu derivate parțiale \\
      de ordinul întîi}}
\end{center}

\vspace{2cm}

\section{Ecuații cu derivate parțiale}

\begin{definition}\label{def:edp1-omogena}
  Fie $ v : U \to \RR^n $ un cîmp de vectori de clasă $ \kal{C}^r(U) $, cu
  $ r \geq 2 $ și $ v_1, \dots, v_n : U \to \RR $ componentele sale.

  Se numește \emph{ecuație cu derivate parțiale de ordinul întîi, liniară și omogenă}
  pentru funcția necunoscută $ u : U \to \RR $, de clasă $ \kal{C}^1 $ egalitatea:
  \begin{equation} \label{eq:edp1-omogena}
    \sum_{i = 1}^n v_i(x_1, \dots, x_n) \frac{\partial u}{\partial x_i}(x_1, \dots, x_n) = 0,
  \end{equation}
  unde $ x = (x_1, \dots, x_n) \in U $.

  Funcția $ u : U \to \RR $ de clasă $ \kal{C}^1 $ care verifică relația de mai sus pentru
  orice $ x \in U $ se numește \emph{soluție} a ecuației pe domeniul $ U $.
\end{definition}

Relevanța integralelor prime pentru ecuațiile cu derivate parțiale rezultă din următoarea.

\begin{theorem}\label{thm:int-prim-edp}
  Orice integrală primă pe $ U $ a sistemului autonom $ x' = v(x) $ este soluție pe $ U $
  a ecuației ~\eqref{eq:edp1-omogena} și, reciproc, orice soluție pe $ U $ a ecuației
  este integrală primă pe $ U $ a sistemului $ x' = v(x) $.
\end{theorem}

În contextul și cu notațiile de mai sus, sistemul diferențial autonom:
\[
  \frac{dx_1}{v_1(x_1, \dots, x_n)} = \frac{dx_2}{v_2(x_1, \dots, x_n)} = \dots = %
  \frac{dx_n}{v_n(x_1, \dots, x_n)}
\]
se numește \emph{sistemul caracteristic} asociat ecuației ~\eqref{eq:edp1-omogena}.

De \textbf{exemplu}, să considerăm ecuația:
\[
  xy \frac{\partial u}{\partial x} - y \sqrt{1 - y^2} \frac{\partial u}{\partial y} + %
  (z \sqrt{1 - y^2} - axy) \frac{\partial u}{\partial z} = 0.
\]

\emph{Soluție:} Sistemul autonom caracteristic este:
\[
  \frac{dx}{xy} = \frac{dy}{-y\sqrt{1 - y^2}} = \frac{dz}{z\sqrt{1 - y^2} - axy}.
\]

Lucrînd cu primele două rapoarte, obținem:
\[
  \frac{dx}{x} = \frac{dy}{-\sqrt{1 - y^2}} \Rightarrow \ln |x| + \arcsin y = c \Rightarrow %
  xe^{\arcsin y} = c_1.
\]

Lucrînd cu ultimele două egalități și folosind soluția pentru $ x $, avem:
\[
  \frac{dy}{-y\sqrt{1 - y^2}} = \frac{dz}{z \sqrt{1 - y^2} - axy} \Rightarrow %
  \frac{dz}{dy} - \frac{z}{y} = ac_1 \frac{e^{-\arcsin y}}{\sqrt{1 - y^2}}.
\]

Aceasta este o ecuație diferențială liniară privind funcția $ z = z(y) $, deci o putem rezolva
cu metoda specifică și găsim soluția generală:
\[
  z = \frac{c_2}{y} \cdot \frac{ac_1 e^{-\arcsin y}}{2y} (y + \sqrt{1 - y^2}) \Leftrightarrow %
  2yz + ax(y + \sqrt{1 - y^2}) = 2c_2.
\]

Soluția generală se alcătuiește, atunci:
\[
  u(x, y, z) = \Phi(xe^{\arcsin y}, 2yz + ax(y + \sqrt{1 - y^2})), \quad \Phi \in \kal{C}^1(U).
\]

Un alt caz de interes este:
\begin{definition}\label{def:edp1-qlin}
  Se numește \emph{ecuație cu derivate parțiale de ordinul întîi cvasiliniară} egalitatea:
  \begin{equation} \label{eq:edp1-qlin}
    \sum_{i = 1}^n g_i(x_1, \dots, x_n, u) \frac{\partial u}{\partial x_i} = g(x_1, \dots, x_n, u),
  \end{equation}
  unde $ g_i, g : D \to \RR $ sînt funcții de clasă $ \kal{C}^1(D \seq \RR^{n+1}) $.
\end{definition}

\begin{definition}\label{def:sol-edp1-qlin}
  Se numește \emph{soluție} a ecuației (\ref{eq:edp1-qlin}) orice funcție de clasă $ \kal{C}^1 $
  definită pe un domeniu $ U \seq \RR^n $, $ u : U \to \RR $, astfel încît $ \forall x \in U $
  să avem $ (x, u(x)) \in D$ și:
  \[
    \sum_{i = 1}^n g_i(x, u(x)) \frac{\partial u}{\partial x_i} (x) = g(x, u(x)), \forall x \in U.
  \]
\end{definition}

Pentru rezolvarea ecuației (\ref{eq:edp1-qlin}) procedăm astfel: căutăm soluția $ u $ sub formă
implicită, $ F(x_1, \dots, x_n, u) = 0 $, unde $F : D \to \RR $ este o funcție de clasă $ \kal{C}^1 $
și $ \frac{\partial F}{\partial u} \neq 0 $ pe $D$. Atunci obținem:
\[
  \frac{\partial u}{\partial x_i} = - \frac{\frac{\partial F}{\partial x_i}}{\frac{\partial F}{\partial u}},
\]
pentru orice $ i = 1, \dots n $, iar ecuația (\ref{eq:edp1-qlin}) devine:
\[
  \sum_{i = 1}^n g_i(x_1, \dots, x_n, u) \frac{\partial u}{x_i} (x_1, \dots, x_n, u) + %
  g(x_1, \dots, x_n, u) \frac{\partial F}{\partial u} (x_1, \dots x_n, u) = 0,
\]
care este o ecuație cu derivate parțiale de ordinul întîi, liniară și omogenă, așadar redusă la
cazul anterior.

De \textbf{exemplu}:
\[
  (1 + \sqrt{z - x - y}) \frac{\partial z}{\partial x} + \frac{\partial z}{\partial y} = 2.
\]

\emph{Soluție:} Scriem sistemul caracteristic:
\[
  \frac{dx}{1 + \sqrt{z - x - y}} = \frac{dy}{1} = \frac{dz}{2}.
\]
Din ultimele două rapoarte, obținem o integrală primă de forma $ z - 2y = c_1 $.

De asemenea, prelucrînd rapoartele, obținem:
\[
  dy = \frac{dz - dx - dy}{- \sqrt{z - x - y}},
\]
de unde rezultă că $ y + 2\sqrt{z - x - y} = c_2 $, care este o altă integrală primă.

Așadar, soluția generală se poate scrie sub forma:
\[
  u(x, y, z) = \Phi(z - 2y, y + 2\sqrt{z - x - y}) = 0.
\]

%%%%%%%%%%%%%%%%%%%%%%%%%%%%%%%%%%%%%%%%%%%%%%%%%%%%%%%%%%%%%%%%%%%%%%

\section{Linii și suprafețe de cîmp}

Putem acum să descriem noțiuni precum liniile de cîmp și liniile de forță, în contextul
ecuațiilor cu derivate parțiale.

\begin{definition}\label{def:linii-cimp}
  Fie $ U \seq \RR^3 $ un domeniu și $ \vec{v} : U \to \RR $ un cîmp de clasă $ \kal{C}^1 $
  fără puncte singulare pe $ U $, adică $ \vec{v} $ nu se anulează pe $ U $), unde:
  \[
    \vec{v} = P(x, y, z) \vec{i} + Q(x, y, z) \vec{j} + R(x, y, z) \vec{k}.
  \]

  Integralele prime (orbitele) care rezultă din sistemul diferențial autonom:
  \begin{equation} \label{eq:sist-cimp}
    \frac{dx}{P(x, y, z)} = \frac{dy}{Q(x, y, z)} = \frac{dz}{R(x, y, z)}, \quad (x, y, z) \in U
  \end{equation}
  se numesc \emph{linii de cîmp} pentru $ \vec{v} $.
\end{definition}

Liniile de cîmp pentru un cîmp vectorial $ \vec{v} $ sînt suporturi de curbe $ \gamma: x = x(t), t \in I$,
în lungul cărora vectorul tangent în fiecare punct $ t $ coincide cu $ \vec{v}(x(t)) $.

Dacă lucrăm cu aplicații fizice concrete, de exemplu, cazul cîmpului electromagnetic, liniile
de cîmp se mai numesc și \emph{linii de forță}.

De \textbf{exemplu}, să determinăm liniile de cîmp pentru cîmpul vectorial:
\[
  \vec{v} = \frac{x^2 + y^2 + z^2}{x^2 z} \vec{i} - \frac{2y}{xz} \vec{j} + \frac{-x^2 + y^2 - z^2}{xz^2}.
\]

\emph{Soluție:} Scriem sistemul simetric asociat:
\[
  \frac{x^2 z dx}{x^2 + y^2 + z^2} = \frac{xzdy}{-2y} = \frac{xz^2 dz}{-x^2 + y^2 - z^2},
\]
pe care îl putem scrie și:
\[
  \frac{2xdx}{x^2 + y^2 + z^2} = \frac{2ydy}{-2y^2} = \frac{2zdz}{-x^2 + y^2 - z^2} = %
  \frac{2xdx + 2ydy + 2zdz}{x^2 + y^2 + z^2 - 2x - x^2 + y^2 - z^2} = %
  \frac{2xdx + 2ydy + 2zdz}{0}
\]

Rezultă că avem o integrală primă de forma $ x^2 + y^2 + z^2 = c_1 $.

Introducînd această integrală primă în primul raport, obținem, împreună cu al doilea:
\[
  \frac{2xdx}{c_1} = -\frac{dy}{y} \Rightarrow y = c_2 e^{-\frac{x^2}{c_1^2}}.
\]

Obținem liniile de cîmp:
\[
  \begin{cases}
    x^2 + y^2 + z^2 &= c_1 \\
    y &= c_2 e^{-\frac{x^2}{c_1^2}}
  \end{cases}
\]

\vspace{1cm}

Trecem acum în dimensiune superioară, introducînd \emph{suprafețele de cîmp}.

\begin{definition}\label{def:supr-cimp}
  O suprafață $ S \seq U $ de clasă $ \kal{C}^1 $, fără puncte singulare se numește
  \emph{suprafață de cîmp} pentru vectorul $ \vec{v} $ dacă în orice punct $ P \in S $,
  vectorul $ \vec{v}(M) $ este tangent la suprafață.
\end{definition}

Pentru a obține ecuația suprafețelor de cîmp procedăm astfel. Presupunem că avem o suprafață $ S $
cu ecuația $ F(x, y, z) = 0 $ în $U$, unde $ F : U \to \RR $ este o funcție de clasă $ \kal{C}^1 $.
Deoarece $ S $ nu are puncte singulare, admite tangentă și normală în orice punct $ M \in S $, iar
un vector director al normalei în $ P $ este dat de ecuația:
\[
  \nabla_M F = \frac{\partial F}{\partial x}(M) \vec{i} + \frac{\partial F}{\partial y}(M)\vec{j} + %
  \frac{\partial F}{\partial z}(M) \vec{k} \neq 0.
\]

Rezultă că $ \vec{v}(M) $ este tangent la $ S $ dacă și numai dacă $ \vec{v}(M) $ este perpendicular
pe vectorul $ \nabla_M F $, adică:
\begin{equation} \label{eq:supr-cimp}
  \vec{v}(M) \cdot \nabla_M F = P \cdot F_x + Q \cdot F_y + R \cdot F_z = 0,
\end{equation}
pentru orice punct $ M = (x, y, z) $. Rezultă că $ S $ este o suprafață de cîmp pentru $ \vec{v} $
dacă și numai dacă este dată de o ecuație de forma $ F(x, y, z) = 0 $, unde funcția $ F : U \to \RR $
este de clasă $ \kal{C}^1 $, are $ \nabla F \neq 0 $ și satisface ecuația (\ref{eq:supr-cimp}),
care se numește \emph{ecuația suprafețelor de cîmp} ale lui $ \vec{v} $. Sistemul diferențial autonom
și simetric (\ref{eq:sist-cimp}) se numește \emph{sistem caracteristic} ale ecuației
(\ref{eq:supr-cimp}), iar liniile de cîmp se mai numesc \emph{curbe caracteristice}.

Determinarea unei suprafețe de cîmp care trece printr-o curbă dată este, de fapt, o problemă
Cauchy pentru ecuația (\ref{eq:supr-cimp}).

\vspace{1cm}

\textbf{Exemplu:} Determinăm soluția ecuației:
\[
  x \frac{\partial z}{\partial x} + y \frac{\partial z}{\partial y} = z,
\]
care trece prin curba $ x^2 + y^2 = 1, z = 2 $.

\emph{Soluție:} Sistemul caracteristic asociat este simplu:
\[
  \frac{dx}{x} = \frac{dy}{y} = \frac{dz}{z}.
\]
Integralele prime se obțin egalînd primele două și ultimele două rapoarte, ca fiind:
\[
  \frac{x}{y} = c_1, \quad \frac{x}{z} = c_2.
\]
Folosindu-le, împreună cu curbele date, avem:
\[
  \begin{cases}
    x &= c_1 y \\
    x &= c_2 z \\
    x^2 + y^2 &= 1 \\
    z &= 2
  \end{cases} \Rightarrow c_1^2 c_2^2 + c_2^2 = %
  \frac{c_1^2}{4} \Rightarrow x^2 + y^2 = \frac{z^2}{4},
\]
care este un con cu vîrful în origine.

% \vspace{1cm}

% \textbf{Exemplu 2:} Fie cîmpul vectorial:
% \[
%   \vec{v} = (x + y) \vec{i} + (y - x) \vec{j} - 2z\vec{k}.
% \]
% Să se determine:
% \begin{enumerate}[(a)]
% \item liniile de cîmp;
% \item linia de cîmp ce conține punctul $ M(1, 0, 1) $;
% \item suprafața de cîmp;
% \item suprafața de cîmp ce conține dreapta $ z = 1, y - x\sqrt{3} = 0 $.
% \end{enumerate}

% \emph{Soluție:} (a) Sistemul caracteristic asociat este:
% \[
%   \frac{dx}{x + y} = \frac{dy}{y - x} = \frac{dz}{-2z} \Rightarrow %
%   \frac{xdx + ydy}{x^2 + y^2} = \frac{dz}{-2z}.
% \]

% Obținem de aici:
% \[
%   \frac{1}{2} \cdot \frac{d(x^2 + y^2)}{x^2 + y^2} = -\frac{1}{2} \frac{dz}{z} \Rightarrow %
%   (x^2 + y^2)z = c_1.
% \]

% Din primele două rapoarte, obținem: $ \frac{dx}{dy} = \frac{x + y}{y - x} $, care este o
% ecuație diferențială omogenă, pentru care facem substituția $ y = tx $. Rezultă:
% \[
%   x \frac{dt}{dx} + t = \frac{t - 1}{t + 1} \Rightarrow \frac{dx}{x} = %
%   -\frac{t + 1}{t^2 + 1} dt,
% \]
% care are soluția $ \ln(x^2 + y^2) + 2 \arctan \frac{y}{x} = c_2 $, care este cea de-a
% doua integrală primă.

% Așadar, liniile de cîmp sînt:
% \[
%   \begin{cases}
%     (x^2 + y^2) z &= c_1 \\
%     \ln(x^2 + y^2) + 2 \arctan \frac{y}{x} &= c_2
%   \end{cases}
% \]

% (b) Cum $ M(1, 0, 1) $ aparține liniei de cîmp, putem afla constantele din sistemul de mai
% sus. Găsim $ c_1 = c_2 = 0 $. \\

% (c) Ecuația suprafeței de cîmp se obține folosind liniile de cîmp:
% \[
%   \Phi \big((x^2 + y^2) z, \ln(x^2 + y^2) + 2 \arctan \frac{y}{x} \big) = 0.
% \]

% (d) Suprafața de cîmp ce conține dreapta dată rezultă cu metoda substituției din
% sistemul:
% \[
%   \begin{cases}
%     (x^2 + y^2)z &= c_1 \\
%     \ln(x^2 + y^2) + 2 \arctan \frac{y}{x} &= c_2 \\
%     z &= 1 \\
%     y - x\sqrt{3} &= 0
%   \end{cases}
% \]

% Obținem:
% \[
%   \ln c_1 + 2\arctan \sqrt{3} = c_2 \Rightarrow z = e^{\arctan \frac{y}{x} - \frac{\pi}{3}},
% \]
% care este ecuația suprafeței de cîmp ce trece prin dreapta dată.

%%%%%%%%%%%%%%%%%%%%%%%%%%%%%%%%%%%%%%%%%%%%%%%%%%%%%%%%%%%%%%%%%%%%%%

\section{Exerciții}

\indent 1. Rezolvați ecuațiile:
\begin{enumerate}[(a)]
\item $ (z - y)^2 \frac{\partial z}{\partial x} + xz \frac{\partial z}{\partial y} = xy $;
\item $ x(y^3 - 2x^3) \frac{\partial z}{\partial x} + %
  y(2y^3 - x^3) \frac{\partial z}{\partial y} = 9z(x^3 - y^3) $;
\item $ 2xu \frac{\partial u}{\partial x} + 2yu \frac{\partial u}{\partial y} = %
  u^2 - x^2 - y^2 $.
\end{enumerate}

\vspace{1cm}

2. Determinați suprafețele de cîmp pentru cîmpurile vectoriale:
\begin{enumerate}[(a)]
\item $ \vec{v} = (x + y + z)\vec{i} + (x - y)\vec{j} + (y - x)\vec{k} $;
\item $ \vec{v} = (x - y)\vec{i} + (x + y) \vec{j} + \vec{k} $;
\item $ \vec{v} = xy^2 \vec{i} + x^2y \vec{j} + z(x^2 + y^2) \vec{k}$.
\end{enumerate}

\vspace{1cm}

3. Determinați suprafața de cîmp a cîmpului vectorial:
\[
  \vec{v} = yz \vec{i} + xz \vec{j} + xy \vec{k},
\]
care trece prin curba de ecuație $ x^2 + y^2 = 4, y = 0 $.

\vspace{1cm}

4. Determinați suprafața de cîmp a cîmpului vectorial:
\[
  \vec{v} = 2zx\vec{i} + 2zy\vec{j} + (z^2 - x^2 - y^2) \vec{k},
\]
care conține cercul de ecuații:
\[
  \begin{cases}
    z &= 0 \\
    x^2 + y^2 - 2x &= 0.
  \end{cases}
\]

\vspace{1cm}

5. Fie cîmpul vectorial:
\[
  \vec{v} = (x + y) \vec{i} + (y - x) \vec{j} - 2z\vec{k}.
\]
Să se determine:
\begin{enumerate}[(a)]
\item liniile de cîmp;
\item linia de cîmp ce conține punctul $ M(1, 0, 1) $;
\item suprafața de cîmp;
\item suprafața de cîmp ce conține dreapta $ z = 1, y - x\sqrt{3} = 0 $.
\end{enumerate}

\emph{Indicație:} Pentru liniile de cîmp, scriem sistemul autonom asociat:
\[
  \frac{dx}{x + y} = \frac{dy}{y-x} = \frac{dz}{-2z}.
\]

Amplificînd prima fracție cu $ x $ și pe a doua cu $ y $, prin adunare, obținem:
\[
  \frac{xdx + ydy}{x^2 + y^2} = \frac{dz}{-2z} \Rightarrow %
  \frac{1}{2} \cdot \frac{d(x^2 + y^2)}{x^2 + y^2} = -\frac{1}{2} \frac{dz}{z}.
\]
Rezultă $(x^2 + y^2)z = c_1$ este o integrală primă.

Din primele două rapoarte, obținem apoi o ecuație diferențială de ordinul întîi,
omogenă:
\[
  \frac{dx}{dy} = \frac{x + y}{y - x},
\]
pe care o putem rezolva cu substituția $y = tx$.

\end{document}